% 编译方式：使用支持中文的 XeLaTeX 连续编译两次。
% 排版依据：旧版论文草稿的文档类、页面设置、标题层级、编号和三线表结构。
% 当前为阶段稿：摘要仅覆盖已完成的问题一；正文分析覆盖全题。
% 问题二至四尚未完成求解，相关方法为研究方案，不代表已取得结果。
% 后续正文在参考文献之前续写，保留现有格式；全题完成后统一修订题目和摘要。
\documentclass[UTF8,a4paper,12pt,fontset=fandol]{ctexart}
\usepackage[left=2.5cm,right=2.5cm,top=2.4cm,bottom=2.4cm]{geometry}
\usepackage{amsmath,amssymb,bm}
\usepackage{booktabs,longtable,tabularx,array,multirow}
\usepackage{graphicx,float,subcaption}
\usepackage{xcolor}
\usepackage{enumitem}
\usepackage{setspace}
\usepackage{fancyhdr}
\usepackage{hyperref}
\usepackage{listings}
\usepackage{caption}
\usepackage{chngcntr}

\hypersetup{colorlinks=true,linkcolor=black,citecolor=black,urlcolor=blue}
\graphicspath{{figures/}{../figures/}}
\captionsetup{font=small,labelsep=quad}
\captionsetup[table]{position=top}
\captionsetup[figure]{position=bottom}
\setstretch{1.22}
\setlength{\parindent}{2em}
\setlength{\parskip}{0pt}
\setlist{nosep,leftmargin=2em}
\pagestyle{fancy}
\fancyhf{}
\cfoot{\thepage}
\renewcommand{\headrulewidth}{0pt}
\renewcommand{\arraystretch}{1.18}
\newcolumntype{Y}{>{\centering\arraybackslash}X}
\newcolumntype{L}{>{\raggedright\arraybackslash}X}
\newcommand{\dif}{\mathop{}\!\mathrm{d}}
\newcommand{\VaR}{\operatorname{VaR}}
\newcommand{\CVaR}{\operatorname{CVaR}}
\numberwithin{equation}{section}
\counterwithin{figure}{section}
\counterwithin{table}{section}
\renewcommand{\thefigure}{\thesection-\arabic{figure}}
\renewcommand{\thetable}{\thesection-\arabic{table}}

\lstset{
  language=Python,
  basicstyle=\ttfamily\scriptsize,
  keywordstyle=\color{blue!70!black},
  commentstyle=\color{green!45!black},
  stringstyle=\color{orange!70!black},
  numbers=left,
  numberstyle=\tiny\color{gray},
  frame=single,
  breaklines=true,
  showstringspaces=false,
  tabsize=4,
  columns=fullflexible
}

\title{\heiti\bfseries 微网与外部电网电力调控策略优化的的研究}
\author{}
\date{}

\begin{document}
\maketitle
\vspace{-2.2em}

% 摘要采用局部紧凑排版；其余正文保留旧稿的默认字号和行距。
% 摘要的一页要求需在完整编译后的页面中最终核验。
\begingroup
\small
\setstretch{1.08}
\begin{abstract}
本文针对光伏发电高峰与居民用电高峰在单日内时间上不一致，以及分时电价下微网购电成本变化的问题，以题目给出的各时段居民用电负荷、光伏发电量预测值和分时电价数据为基础，先计算各时段的电力供需差额，再建立储能优化调度模型。模型以外部电网购电量和储能充放电量为决策变量，以全天购电费用最小为目标，并设置电力供需平衡、储能容量、充放电功率和剩余电量等约束，求得各时段的购电与储能运行方案。在此基础上，将无储能方案作为对照，通过约束回代和充放电效率敏感性分析，检验模型结果的可行性和经济性，并分析储能效率变化时调度方案的调整情况。

首先，将全天划分为144个十分钟时段，统一功率与电量的换算口径，以购电量、充电量、放电量、弃光量及储能状态为变量，建立\textbf{混合整数线性规划模型}。模型引入二元变量约束充放电互斥，综合考虑供需平衡、储能状态递推、容量边界、充放电功率上限及日初日末储能相等条件。在允许弃光、不向外部电网售电的假设下，将充电和放电单程效率均设为90\%，储能运行区间设为1200--10800千瓦时，最大充放电功率均为5000千瓦，确保调度方案满足设备运行条件。

其次，采用\textbf{两层词典序优化方法}：第一层最小化全天购电费用；第二层在购电费用不超过第一层最优值加给定容差的条件下，最小化储能充放电总吞吐量，避免人为设定折旧权重改变首要目标。模型通过低价时段购电充能、光伏富余时段吸纳余电及高价时段放电，实现跨时段能量配置，并生成逐时段购电方案与分时段充放电汇总。

求解结果表明，优化后全天购电费用为\textbf{35126.95元}，较无储能基准48052.05元降低12925.10元，\textbf{降幅为26.90\%}；全天购电量为59482.70千瓦时，充电量与放电量分别为20740.67千瓦时和16799.94千瓦时。在给定光伏预测情景下，全天弃光量为零，日末储能恢复至6000千瓦时，全部时段均满足负荷需求，且不存在同时充放电现象。

最后，对输出方案进行数值检验，供需平衡与储能递推的最大残差分别约为$2.27\times10^{-13}$和$8.31\times10^{-13}$千瓦时，储能状态全程满足上下界约束。效率敏感性分析显示，单程效率由80\%提高至95\%时，最优购电费用由38223.65元降至33767.03元；若将题设90\%解释为往返效率，费用为33801.50元，表明效率口径会影响经济性评估。研究表明，该模型能够在满足储能物理约束的条件下有效降低购电成本、促进光伏消纳，具有结构清晰、结果可解释及便于复现的优势，可为后续考虑预测误差与波动电价的调度研究提供基础。

\par\medskip
\noindent\textbf{关键词：}混合整数线性规划；两层词典序优化；微网电力调控；储能调度；光伏消纳
\end{abstract}
\endgroup
\clearpage

\section{问题重述}
\label{sec:restatement}

\subsection{问题背景}

含光伏与储能的并网型微网能够在本地协调发电、用电及能量存储，但其经济运行受到供需错配与预测不确定性的共同制约。一方面，光伏出力具有明显的日内变化，发电高峰未必与居民用电高峰重合，局部时段可能出现电能富余，而其他时段仍需依赖外部电网。另一方面，储能虽然能够转移电能的使用时刻，却受到容量、充放电功率和能量转换损耗的限制。因此，仅依据“低价充电、高价放电”的经验规则，难以协调光伏消纳、购电支出和后续时段的供电需求。

当调度由确定性单日计划扩展至连续运行时，问题进一步表现为预测与决策之间的相互影响。计划购电偏少会增加高价紧急购电的需求；计划偏多则可能形成无法充分利用、但仍需支付费用的购电量。即使日内能够更新光伏预报，调整计划也伴随着额外交易费用，因此，预测更新能否带来实际收益，需要通过完整的调度成本衡量。

近年来的相关研究为上述问题提供了方法依据。2022年发表于《可再生能源》的研究将实际负荷、可再生能源数据及预测误差纳入经济模型预测控制框架，强调利用滚动决策适应实际运行状态\cite{economic2022}。2023年的风险感知微网能源管理研究指出，应系统考虑发电与用电预测的不确定性，同时控制随机优化的计算负担\cite{riskaware2023}。据此，本题需要围绕“物理约束下的经济调度---预测误差下的补救决策---信息更新下的动态调整”建立统一研究框架。

\subsection{研究对象与约束条件}

研究对象为由小区负荷、光伏发电、储能设备及外部电网构成的并网型微网。调度以十分钟为基本时间间隔，每日包含144个时段，需要联合确定各时段的计划购电量和储能充放电量，并在需要时安排计划调整与紧急购电。

储能设备额定容量为12000千瓦时，允许运行电量为1200--10800千瓦时，最大充电功率和最大放电功率均为5000千瓦。2025年1月1日零时的初始储电量为6000千瓦时，题设充放电效率为90\%。其中，问题一还要求日初与日末储电量相同。

数据包括典型日的负荷、光伏预测与电价，以及2025年全年的实际负荷、实际光伏出力、多时点光伏预报和波动电价。全年策略需按要求输出2月1日至12月31日的调度记录，并展示3月20日、6月21日、9月23日和12月21日的代表性结果。预测数据与实际数据应按发布时间及目标时段分别管理，避免把事后观测值作为事前已知条件。

\subsection{核心研究任务}

\textbf{研究任务一：确定性条件下的购电与储能协同优化，对应原题问题一。}

在当日负荷、电价及光伏预测已知的条件下，建立满足供需平衡与储能运行限制的单日调度模型。重点研究储能应在何时吸收富余光伏、何时利用低价电补能，以及何时释放电能替代高价购电，使日初日末储电量相同条件下的全天购电费用最小，并形成可执行的逐时段调度方案。

\textbf{研究任务二：预测不确定性下的日前计划与日内调整，对应原题问题二、三。}

首先，在每日电价曲线相同、负荷与光伏实际出力变化的条件下，利用历史信息制定零时购电计划，并以交易时刻电价五倍的紧急购电弥补实际供电缺口。其次，引入零时、六时、十二时和十八时发布的光伏预报，在已执行决策不可更改的前提下更新剩余时段的购电与储能安排。研究重点是平衡提前采购、日内调整和紧急补购之间的成本关系，并判断增加预报时点的边际经济价值。

\textbf{研究任务三：波动电价下的策略扩展与适应性评价，对应原题问题四。}

将固定日内电价曲线替换为实际波动电价，分别研究仅进行日前计划和允许日内调整两种机制。重点分析电价不确定性与供需预测误差对储能使用时机的共同影响，在一致的数据范围、运行边界和结算口径下比较两种策略，为不同交易条件下的购电决策提供依据。

% 问题分析采用独立页面与局部紧凑排版；不混入尚未完成的求解结果。
\clearpage
\section{问题分析}
\label{sec:analysis}
\begingroup
\small
\setstretch{1.08}

\textbf{总体思路。}本题整体采用“数据处理与因果预测---约束优化---滚动执行---对照检验”的思路。以供需平衡和储能状态递推构成统一物理模型，逐步增加紧急购电、计划调整和电价不确定性。购电、调整与补救费用共同构成总成本，属于具有多个成本分项的单目标优化；问题一再以两层词典序方式处理等价方案选择。

\textbf{问题1分析。}关键矛盾是降低当前购电费用与保留后续可用储能之间的跨时段权衡。首先统一功率、电量与时间索引，识别净负荷及光伏富余区间；其次建立购电、充放电和储能状态变量，以二元变量保证充放电互斥，并施加容量、功率及首末储能约束；最后先最小化购电费，再在规定成本容差内最小化吞吐量，通过约束回代和线性松弛下界检验方案。

\textbf{问题2分析。}关键矛盾是计划不足引发高价补购，而过度采购仍需承担计划费用。先按时间顺序利用历史数据预测负荷和光伏，再根据预测残差构造保留时间相关性的情景；随后建立日前计划与实时补救相衔接的两阶段模型。执行时仅依据当时可用信息调整储能和紧急购电，保持跨日状态连续，避免利用未来实际值获得不可实现的策略优势。

\textbf{问题3分析。}关键矛盾是新预报可能降低供需偏差，但修改计划需要支付费用。先将整点预报统一至十分钟网格，再于各发布时间更新剩余时域的预测，固定已执行决策，以当前储能和有效购电计划为起点进行滚动优化。结算时区分原始计划、增减调整量和紧急购电。新增预报时点通过成对回测比较净成本变化，其收益须覆盖信息获取及调整成本。

\textbf{问题4分析。}关键矛盾是储能同时承担电价套利与供电缓冲功能，单纯追逐价差可能削弱应对缺口的能力。先明确未来电价的可获得范围，对未知电价建立预测及联合误差情景；再分别嵌入问题二、三的交易结构，保持一致的评价条件。最后结合总成本、紧急购电、尾部成本及参数敏感性，分析策略对价格波动的适应能力。

% 以下流程框架默认不排印，避免与分析正文重复并占用单页篇幅。
% 如后续需要流程图，可据此绘制并通过图环境插入。
% 问题一：数据与单位核验→确定性混合整数优化→词典序择优→约束及下界检验。
% 问题二：历史信息→预测与误差情景→日前计划→实时补救→跨日状态更新。
% 问题三：新预报到达→固定已执行决策→剩余时域优化→费用结算→新增预报价值评价。
% 问题四：电价信息边界判断→联合预测与情景生成→分别运行两种交易机制→统一对照评价。
\par
\endgroup
\clearpage

\section{模型假设}
\label{sec:assumptions}

以下五条用于补充题目未充分明确的建模边界。容量、功率限制及供电要求属于题设约束，不再作为假设重复列出。

\begin{enumerate}[label=\textbf{假设\arabic*：},leftmargin=4.5em]
\item \textbf{采用十分钟离散调度，所有决策遵守信息可得性。}

\textbf{假设内容：}每个十分钟时段内，以给定功率或预测功率代表该时段的平均运行水平。问题一将给定光伏预测作为确定性输入；问题二至四的决策仅使用决策时点及以前的观测和已发布预报。问题四中未明确提前公布的未来电价按未知量处理。整点光伏预报通过固定的时间转换规则映射至十分钟网格。

\textbf{合理性依据：}题目要求以十分钟粒度输出购电方案，而光伏预报具有明确的发布时间。统一时间尺度是连接预测与调度的必要条件，限制信息范围则符合实际决策顺序。

\textbf{对模型的影响：}将连续运行过程转化为有限维优化问题，并防止未来信息泄漏。模型不描述十分钟以内的瞬时波动；预报插值也不意味着获得了额外真实信息，其误差应纳入后续回测。

\item \textbf{储能采用恒定效率模型，主方案将90\%解释为充、放电各自的单程效率。}

\textbf{假设内容：}充电效率和放电效率均取0.9，忽略自放电及效率随温度、功率和储能状态的变化。研究期内可用容量保持不变，不直接引入缺少参数依据的货币化衰减费用；问题一以吞吐量作为等价方案的次级筛选指标。

\textbf{合理性依据：}题目仅提供一个效率参数，未提供效率曲线、自放电率和循环寿命数据，因此恒定参数模型具有明确的数据基础。考虑到“充放电效率”存在解释空间，另设往返效率为90\%的敏感性方案，此时充、放电效率均取其平方根。

\textbf{对模型的影响：}储能递推保持线性，便于精确求解和检查能量守恒。但该模型不能直接预测电池寿命，也不能将较低吞吐量等同于已验证的寿命延长。

\item \textbf{微网按单节点能量系统描述，允许弃光及处置无法利用的剩余电量，不向外网售电。}

\textbf{假设内容：}忽略微网内部线路损耗、电压与无功功率限制；外网在研究情景内可正常供电，且不另设题目未给出的购电功率上限。光伏富余时允许弃光；已购买但无法用于负荷或储能的电量按未利用电量记录，不产生售电收入。

\textbf{合理性依据：}题目提供的是聚合负荷、光伏和储能数据，没有网络拓扑、线路参数及售电价格。问题二规定普通购电费用按计划量计算，因此预测偏差导致的采购富余不能自动获得退款或收益。

\textbf{对模型的影响：}可以聚焦购电与储能的能量调度，并保证富余电量在平衡方程中具有明确去向。弃光与未利用购电应分别核算，以免夸大光伏消纳能力。结论适用于上述并网条件，不直接覆盖停电或网络拥塞情景。

\item \textbf{计划调整采用统一的结算解释，同一调整量不重复收费。}

\textbf{假设内容：}主方案保留零时原始计划购电费；有效购电量相对原计划上调的部分按交易时刻电价的1.5倍计费，下调的部分按0.5倍计收违约费用，且不另行假定原计划费用返还。紧急购电独立按五倍价格计费。多次滚动优化仅形成候选方案，费用按各执行时段最终生效的调整量核算。

\textbf{合理性依据：}题目列明了计划费用、调整费用及紧急购电费用，但对下调退款和同一未来时段反复修改的累计收费方式未作完整说明。采用统一解释可形成可复现的费用账本，并区分优化过程与实际交易。

\textbf{对模型的影响：}上、下调费用可用非负增减变量线性表示，避免将完整调整后购电量与原始计划重复计价。由于该假设可能影响调整积极性，后续应对“下调退还相应基础电费”等替代口径进行敏感性比较。

\item \textbf{全年运行保持储能跨日连续，设置统一的评价期末条件。}

\textbf{假设内容：}问题一严格满足单日首末储能相等；问题二至四以前一天末的实际储电量作为次日初值，不逐日强制恢复到6000千瓦时。全年模拟自题设初始状态启动，1月用于初始预测训练及状态衔接，规定评价期末储能恢复至题设初始水平，并对所有对照策略采用相同规则。

\textbf{合理性依据：}实际储能不会在午夜自动复位，逐日重置会切断相邻日期间的能量联系；而完全不限制评价期末状态，又可能使模型通过耗尽电池降低账面费用。

\textbf{对模型的影响：}保留跨日调度能力，提高策略比较的公平性。评价期末条件属于额外的比较约定，需检查其对末期调度的影响；滚动时域末端的储能价值应单独处理，不能机械地替换为每日首末相等。
\end{enumerate}

\section{符号说明}
\label{sec:notation}

表\ref{tab:notation}给出全题共用及后续模型拟采用的主要符号。功率与电量、交流侧电量与电池内部电量以及预测值与实际执行值分别表示，避免不同建模阶段出现口径混用。

\begingroup
\small
\setlength{\tabcolsep}{4pt}
\begin{longtable}{>{\raggedright\arraybackslash}p{0.25\textwidth}>{\raggedright\arraybackslash}p{0.53\textwidth}>{\raggedright\arraybackslash}p{0.15\textwidth}}
\caption{主要符号及其含义}\label{tab:notation}\\
\toprule
符号 & 含义 & 单位\\
\midrule
\endfirsthead
\multicolumn{3}{c}{表\thetable\quad 主要符号及其含义（续）}\\
\toprule
符号 & 含义 & 单位\\
\midrule
\endhead
\midrule
\multicolumn{3}{r}{续下页}\\
\endfoot
\bottomrule
\endlastfoot
$t,\mathcal T$ & 十分钟时段索引及所研究时段集合 & 无量纲\\
$k,\tau_k$ & 滚动决策序号及对应的决策时点 & 无量纲；时刻\\
$\omega,\Omega$ & 不确定性情景索引及情景集合 & 无量纲\\
$\pi_\omega$ & 情景$\omega$的发生概率 & 无量纲\\
$\Delta t$ & 调度时间间隔，取$1/6$小时 & 小时\\
$\mathcal I_k$ & 决策时点$\tau_k$已获得的观测与预报信息集合 & ---\\
$P_t^L,P_t^{PV}$ & 时段平均负荷功率、可用光伏功率 & 千瓦\\
$L_t,R_t$ & 时段负荷电量、可用光伏电量 & 千瓦时\\
$\widehat L_{t\mid k},\widehat R_{t\mid k}$ & 在决策时点$\tau_k$对时段$t$的负荷、光伏电量预测 & 千瓦时\\
$N_t$ & 净负荷电量，定义为$L_t-R_t$ & 千瓦时\\
$p_t,\widehat p_{t\mid k}$ & 时段结算电价及在决策时点形成的对应电价预测 & 元／千瓦时\\
$g_t^0$ & 每日零时确定的时段原始计划购电量 & 千瓦时\\
$g_{t\mid k}$ & 第$k$次决策形成的时段非紧急购电计划量 & 千瓦时\\
$g_t$ & 时段最终生效的非紧急购电量 & 千瓦时\\
$\delta_t^+,\delta_t^-$ & 最终生效购电量相对零时原计划的上调量、下调量 & 千瓦时\\
$e_t$ & 时段紧急购电量 & 千瓦时\\
$c_t,d_t$ & 时段交流侧充电量、交流侧放电量 & 千瓦时\\
$w_t,s_t$ & 时段弃光电量、已购买但未利用的电量 & 千瓦时\\
$E_t$ & 时段结束时的电池内部储电量 & 千瓦时\\
$E^{\mathrm{ini}},E^{\mathrm{ter}}$ & 所研究区间的初始、目标终端储电量 & 千瓦时\\
$E^{\min},E^{\max}$ & 允许储电量下限、上限，分别为1200、10800 & 千瓦时\\
$P^{ch,\max},P^{dis,\max}$ & 最大充电功率、最大放电功率，均为5000 & 千瓦\\
$\eta_{ch},\eta_{dis}$ & 充电效率、放电效率 & 无量纲\\
$u_t^{ch},u_t^{dis}$ & 时段充电、放电状态二元变量 & 无量纲\\
\shortstack[l]{$C_{\mathrm{plan}},C_{\mathrm{adj}}$\\$C_{\mathrm{emg}}$} & 评价区间内计划购电费、调整费、紧急购电费 & 元\\
$C_{\mathrm{tot}}$ & 评价区间内实际总购电费用 & 元\\
$C^*,\varepsilon_C$ & 问题一第一层最优购电费、第二层允许的费用容差 & 元\\
$Q$ & 评价区间内储能交流侧充放电总吞吐量 & 千瓦时\\
$V_\tau$ & 增加预报时点$\tau$所带来的预期净成本节约 & 元\\
\end{longtable}
\endgroup

为保持后续公式一致，统一以下口径：
\begin{enumerate}
\item \textbf{功率与电量严格区分。}$L_t=P_t^L\Delta t$、$R_t=P_t^{PV}\Delta t$。购电及充放电变量均表示时段电量，计算费用时不再重复乘以时间间隔。
\item \textbf{储能侧与交流侧严格区分。}$c_t,d_t$为交流侧电量，$E_t$为电池内部电量，效率仅在储能状态递推中计入。
\item \textbf{计划与执行严格区分。}未调整时$g_t=g_t^0$；允许调整时$g_t=g_t^0+\delta_t^+-\delta_t^-$。两类调整量分别表示同一差额的正、负部分，不应同时为正。
\item \textbf{预测与情景保留索引。}帽号表示预测，竖线后的$k$表示预测形成时点；情景模型中的变量按需增加$\omega$索引，实际执行结果不得混用情景预测值。
\end{enumerate}

% ==================== 后续正文接续位置 ====================
% 用户要求后续步骤继续以本文件格式输出并更新本稿。
% 下一阶段新增章节放在此处，参考文献始终置于完整正文之后。
% 正文一级标题使用 section，二级标题使用 subsection，三级标题使用 subsubsection。
% 公式使用 equation、align 等环境；表格使用三线表；全局符号与上表同步。
% 未完成求解的内容不得写成已完成结果；新增结论必须能够追溯到结果文件。

\clearpage
\begin{thebibliography}{99}
% 以下保留前述已核验的题名中文译名、年份、卷页与原文链接。
% 正式定稿前补齐作者和原文题名，并按最终要求统一参考文献著录格式。
\bibitem{economic2022}
基于实际需求、可再生能源数据及预测误差的孤岛微网经济预测控制（题名中文译名）.
可再生能源, 2022, 194: 647--658.
\href{https://doi.org/10.1016/j.renene.2022.05.103}{原文链接}.

\bibitem{riskaware2023}
含电池储能与可再生能源的微网风险感知随机能源管理（题名中文译名）.
国际自动控制联合会论文集, 2023, 56(2): 8445--8450.
\href{https://doi.org/10.1016/j.ifacol.2023.10.1042}{原文链接}.
\end{thebibliography}

\end{document}